\titlespacing*{\chapter}
{0pt}
{-40pt}
{20pt}


\label{ch:chapter_one}%
\newpage
\chapter{Metodologia di progetto}

\section{Strategia di ottimizzazione}

L’ottimizzazione è stata implementata in MATLAB utilizzando l’algoritmo genetico (\texttt{ga}) per individuare una configurazione iniziale soddisfacente e successivamente è stato usato il solutore \texttt{fmincon} per il raffinamento locale della soluzione. Il codice completo è riportato in Appendice~\ref{app:codici}.

Dal momento che l’algoritmo genetico è un metodo di ottimizzazione stocastico, il risultato ottenuto dipende dalla popolazione iniziale e dalle operazioni casuali di selezione e mutazione. Per ridurre tale dipendenza ciascun problema di ottimizzazione è stato risolto mediante un numero elevato di esecuzioni indipendenti dell’algoritmo (150 esecuzioni). Al termine delle simulazioni è stata scelta la configurazione che ha fornito il valore della funzione obiettivo minore. La verifica della robustezza della soluzione mediante l’andamento del minimo cumulativo è riportata in Appendice~\ref{app:convergenza_turbojet}.

Il turbogetto e il ramjet sono stati ottimizzati mediante due procedure indipendenti. In una prima fase è stato eseguito il dimensionamento del turbogetto nelle condizioni di crociera subsonica, ottenendo le principali caratteristiche geometriche del motore. Successivamente è stato sviluppato il modello del ramjet.

L'unico collegamento tra i due sistemi propulsivi è rappresentato dal raggio esterno della turbina del turbogetto, utilizzato come vincolo geometrico per il dimensionamento del condotto anulare. 


%-----------------------------------------------------
\section{Dimensionamento e ottimizzazione del turbogetto}
\label{sec:ottimizzazione_turbojet}

L'ottimizzazione del turbogetto è stata orientata al contenimento degli
ingombri. La sola minimizzazione del TSFC avrebbe infatti favorito
configurazioni caratterizzate da elevate portate d'aria e maggiori
dimensioni; è stata pertanto massimizzata la spinta specifica, imponendo
il TSFC come vincolo progettuale.

Le variabili di progetto e i relativi intervalli di variazione sono
riportati in Tabella~\ref{tab:var_turbojet}.

\begin{table}[H]
\centering
\fontsize{11pt}{11pt}\selectfont
\renewcommand{\arraystretch}{1}
\caption{Variabili di progetto considerate durante l'ottimizzazione del turbogetto.}
\label{tab:var_turbojet}
\begin{tabular}{lccc}
\toprule
\textbf{Parametro} & \textbf{Simbolo} & \textbf{Intervallo} & \textbf{Unità} \\
\midrule

Portata massica aria
& $\dot{m}$
& $20-60$
& kg/s \\

Rapporto di compressione compressore
& $\pi_c$
& $3-30$
& -- \\

Velocità periferica al tip del compressore
& $U_{tip}$
& $330-410$
& m/s \\

Temperatura in camera di combustione
& $T_{04}$
& $1200-1700$
& K \\

\bottomrule
\end{tabular}
\end{table}

Gli intervalli sono stati affinati mediante analisi preliminari. La portata
$\dot{m}=20{-}60\ \mathrm{kg/s}$ comprende le configurazioni di interesse
e consente di soddisfare il requisito di spinta, influenzando principalmente
il dimensionamento delle sezioni di passaggio. 



Per \(\pi_c\) è stato imposto un limite superiore pari a 30, oltre il quale l'incremento delle prestazioni risulta limitato rispetto all'aumento del numero di stadi e della complessità del compressore.

La temperatura totale in camera di combustione \(T_{04}\) è stata limitata all'intervallo \(1200{-}1700\ \mathrm{K}\). Il limite inferiore garantisce un'adeguata spinta specifica, mentre quello superiore è imposto dai limiti termici della turbina, considerando l'impiego di superleghe monocristalline e sistemi di raffreddamento.


$U_{tip}=330{-}410\ \mathrm{m/s}$ è stato scelto per
consentire il rispetto dei principali vincoli aerodinamici del compressore.

I principali vincoli adottati durante l'ottimizzazione sono riportati in
Tabella~\ref{tab:vincoli_turbojet}.

\begin{table}[H]
\centering
\caption{Principali vincoli adottati durante l'ottimizzazione del turbogetto.}
\fontsize{11pt}{11pt}\selectfont
\renewcommand{\arraystretch}{1}
\label{tab:vincoli_turbojet}
\begin{tabular}{lll}
\toprule
\textbf{Categoria} & \textbf{Parametro} & \textbf{Limite} \\
\midrule

Prestazioni
& TSFC
& $<0.140\ \mathrm{kg/(N\,h)}$ \\

Prestazioni
& Spinta
& $=25.0\ \mathrm{kN}$ \\

Aerodinamica
& $M_{rel}$
& $<1$ \\

Aerodinamica
& De Haller
& $>0.720$ \\

Aerodinamica
& $C_p$
& $<0.5$ \\

Aerodinamica
& $D_{rotore}$
& $<0.6$ \\

Aerodinamica
& $D_{statore}$
& $<0.6$ \\

Combustione
& $t_{res}$
& $>8\ \mathrm{ms}$ \\

Combustione
& $\phi$
& $0.7<\phi<1.1$ \\

Geometria
& $N_{stadi,c}$
& $<15$ \\

Geometria
& $N_{stadi,t}$
& $<5$ \\

Geometria
& $A_{35},A_{5},A_{ref}$
& $<0.8\ \mathrm{m^2}$ \\

Geometria
& $h_3$
& $> 0.03 $ m \\

\bottomrule
\end{tabular}
\end{table}

I vincoli aerodinamici sono stati introdotti per garantire il corretto funzionamento del compressore e della camera di combustione. Inoltre è stata data maggiore importanza ai vincoli geometrici, al fine di contenere gli ingombri complessivi del motore.

I range di Tabella~1.2 sono stati progressivamente ridotti durante l'ottimizzazione, imponendo vincoli più stringenti sull'ingombro del propulsore. È stata privilegiata la dimensione radiale, per il suo impatto sull'integrazione e sulla resistenza aerodinamica delle gondole.

Inoltre i limiti imposti al numero di stadi delle turbomacchine e alle principali aree caratteristiche del flusso sono stati definiti in modo da contenere indirettamente l'ingombro assiale del propulsore senza penalizzare le prestazioni ottenibili.

Il TSFC del turbogetto è stato vincolato a \(0.14\ \mathrm{kg/(N\,h)}\), valore scelto durante l'ottimizzazione come compromesso tra consumo e ingombro. Soglie superiori aumentavano il consumo, mentre soglie inferiori non determinavano una significativa riduzione delle dimensioni del motore.



%-----------------------------------------------------
\section{Dimensionamento e ottimizzazione del ramjet}
L'ottimizzazione del ramjet ha seguito uno sviluppo simile a quella del turbogetto. La missione impone due vincoli: la spinta pari a 65 kN, e il consumo di carburante in volo supersonico che deve rispettare i limiti riportati nella Tabella  \ref{tab:vincoli_ramjet}, il tutto a una quota di $25000$ m.

\begin{table}[H]
\centering
\fontsize{11pt}{11pt}\selectfont
\renewcommand{\arraystretch}{1}
\caption{Vincoli di missione}
\label{tab:vincoli_ramjet}
\begin{tabular}{lll}
\toprule
\textbf{Categoria} & \textbf{Valore} \\
\midrule

Range
& 5000km \\

Consumo configurazione 2 motori ramjet
& 48000kg \\

Mach di volo
& 3.5\\

Spinta
& $65\ \mathrm{kN}$ \\

Quota
& $25000\ \mathrm{m}$\\

\bottomrule
\end{tabular}
\end{table}

Tali condizioni possono essere ridefinite come vincolo sul TSFC del ramjet. Come calcolato in appendice \ref{ap:appendixA} il ramjet deve presentare un TSFC inferiore a 0.277 durante il volo supersonico per garantire il successo della missione. Date tali condizioni, la spinta specifica è stata scelta come funzione obiettivo al fine di minimizzare le dimensioni del motore, rispettando i vincoli sul TSFC. In particolare si sono imposti i seguenti vincoli:

\begin{table}[H]
\centering
\fontsize{11pt}{11pt}\selectfont
\renewcommand{\arraystretch}{1}
\caption{Vincoli adottati durante l'ottimizzazione del ramjet.}
\label{tab:vincoli_ramjet2}
\begin{tabular}{lll}
\toprule
\textbf{Categoria} & \textbf{Parametro} & \textbf{Limite} \\
\midrule

Prestazioni
& TSFC
& $<0.277\ \mathrm{kg/(N\,h)}$ \\

Prestazioni
& Spinta
& $ = 65\ \mathrm{kN}$ \\


Geometria
& $D_{Eflusso}$
& $<1.6\ \mathrm{m}$ \\

\bottomrule
\end{tabular}
\end{table}

Il vincolo geometrico sul diametro di efflusso è stato imposto poiché, ad alte quote un ugello adattato comporta un notevole salto di pressione il quale rende l'area di efflusso molto grande. È stato quindi deciso di introdurre un vincolo geometrico e di far sotto espandere il flusso così da garantire aree più piccole.

I parametri di progetto per l'ottimizzazione sono i 
seguenti:
\begin{table}[H]
\centering
\fontsize{11pt}{11pt}\selectfont
\renewcommand{\arraystretch}{1}
\caption{Variabili di progetto utilizzati per l'ottimizzazione del ramjet.}
\label{tab:var_turbojet}
\begin{tabular}{lccc}
\toprule
\textbf{Parametro} & \textbf{Simbolo} & \textbf{Intervallo} & \textbf{Unità} \\
\midrule

Portata massica aria &
$\dot{m}$ &
$5 - 400$ &
kg/s \\

Temperatura di combustione &
$T_{04}$ &
$1500 - 2100$ &
K \\

Fattore di espansione motore &
$\pi_{esp}$ &
$4 - 10$ & 
-- \\

\bottomrule
\end{tabular}
\end{table}

La portata massica è stata scelta come parametro da ottimizzare per lo stesso motivo del turbogetto, per garantire la spinta richiesta. Non ha alcuna influenza sul ciclo termodinamico e quindi sulla spinta specifica.

La temperatura totale di combustione nel post-bruciatore è stata imposta in un intervallo compreso tra 1500 K e 2100 K al fine di massimizzare la spinta specifica. 

Infine, il fattore di espansione nell'ugello è stato imposto per garantire il vincolo geometrico sul diametro di efflusso. In particolare è un fattore moltiplicativo applicato alla pressione di uscita:
\begin{equation*}
    P_{e} = P_s\pi_{esp}
\end{equation*}
dove $P_s$ è la pressione statica a $25000$ m di quota.

\subsection{Ottimizzazione della presa supersonica}
La presa supersonica adotta una configurazione a compressione esterna su tre rampe, scelta come compromesso tra complessità geometrica e recupero di pressione totale. La ripartizione della deflessione totale tra i tre urti obliqui è stata determinata secondo il criterio di Oswatitsch, secondo cui, a parità di Mach di volo e deflessione complessiva, il recupero di pressione totale è massimizzato quando tutti gli urti hanno la medesima componente di Mach normale $M_n$. Il problema si riduce così alla ricerca di un unico parametro libero, il rapporto di pressione statica associato a $M_n$, che massimizza il recupero di pressione totale complessivo del sistema di tre urti obliqui più urto normale terminale. 

Tale ricerca è stata condotta numericamente mediante una funzione di ottimizzazione scalare vincolata (\texttt{fminbnd}), che esplora l'intervallo di rapporti di pressione statica ammissibili, escludendo agli estremi le soluzioni fisicamente non accettabili (onde non compressive da un lato, distacco dell'urto o flusso a valle subsonico prematuro dall'altro). Ad ogni tentativo la funzione obiettivo ricostruisce in cascata i tre urti obliqui e l'urto normale terminale a partire dal valore di $M_n$ candidato, verificando la realizzabilità fisica della soluzione (angoli reali, Mach a valle supersonico), e restituisce il recupero di pressione totale risultante. Il valore di $M_n$ che massimizza tale recupero viene infine utilizzato per la ricostruzione dettagliata dei tre urti.